\section{Invertible sheaves and morphisms into Proj}
\label{section-invertible-proj}

\noindent
Let $T$ be a scheme and let $\mathcal{L}$ be an invertible sheaf
on $T$. For a section $s \in \Gamma(T, \mathcal{L})$ we denote
$T_s$ the open subset of points where $s$ does not vanish. See
Modules, Lemma \ref{modules-lemma-s-open}. We can view the following
lemma as a slight generalization of Lemma \ref{lemma-apply}.
It also is a generalization of Lemma \ref{lemma-morphism-proj}.

\begin{lemma}
\label{lemma-invertible-map-into-proj}
Let $A$ be a graded ring.
Set $X = \text{Proj}(A)$.
Let $T$ be a scheme.
Let $\mathcal{L}$ be an invertible $\mathcal{O}_T$-module.
Let $\psi : A \to \Gamma_*(T, \mathcal{L})$ be a homomorphism
of graded rings. Set
$$
U(\psi) = \bigcup\nolimits_{f \in A_{+}\text{ homogeneous}} T_{\psi(f)}
$$
The morphism $\psi$ induces a canonical morphism of schemes
$$
r_{\mathcal{L}, \psi} :
U(\psi) \longrightarrow X
$$
together with a map of $\mathbf{Z}$-graded $\mathcal{O}_T$-algebras
$$
\theta :
r_{\mathcal{L}, \psi}^*\left(
\bigoplus\nolimits_{d \in \mathbf{Z}} \mathcal{O}_X(d)
\right)
\longrightarrow
\bigoplus\nolimits_{d \in \mathbf{Z}} \mathcal{L}^{\otimes d}|_{U(\psi)}.
$$
The triple $(U(\psi), r_{\mathcal{L}, \psi}, \theta)$ is
characterized by the following properties:
\begin{enumerate}
\item For $f \in A_{+}$ homogeneous we have
$r_{\mathcal{L}, \psi}^{-1}(D_{+}(f)) = T_{\psi(f)}$.
\item For every $d \geq 0$ the diagram
$$
\xymatrix{
A_d \ar[d]_{(\ref{equation-global-sections})} \ar[r]_{\psi} &
\Gamma(T, \mathcal{L}^{\otimes d}) \ar[d]^{restrict} \\
\Gamma(X, \mathcal{O}_X(d)) \ar[r]^{\theta} &
\Gamma(U(\psi), \mathcal{L}^{\otimes d})
}
$$
is commutative.
\end{enumerate}
Moreover, for any $d \geq 1$ and any open subscheme $V \subset T$
such that the sections in $\psi(A_d)$ generate $\mathcal{L}^{\otimes d}|_V$
the morphism $r_{\mathcal{L}, \psi}|_V$ agrees with the morphism
$\varphi : V \to \text{Proj}(A)$ and the map $\theta|_V$ agrees with the map
$\alpha : \varphi^*\mathcal{O}_X(d) \to \mathcal{L}^{\otimes d}|_V$
where $(\varphi, \alpha)$ is the pair
of Lemma \ref{lemma-converse-construction}
associated to
$\psi|_{A^{(d)}} : A^{(d)} \to \Gamma_*(V, \mathcal{L}^{\otimes d})$.
\end{lemma}

\begin{proof}
Suppose that we have two triples $(U, r : U \to X, \theta)$
and $(U', r' : U' \to X, \theta')$ satisfying (1) and (2).
Property (1) implies that $U = U' = U(\psi)$ and that
$r = r'$ as maps of underlying topological
spaces, since the opens $D_{+}(f)$ form a basis for the topology
on $X$, and since $X$ is a sober topological space (see
Algebra, Section \ref{algebra-section-proj}
and
Schemes, Lemma \ref{schemes-lemma-scheme-sober}).
Let $f \in A_{+}$ be homogeneous. Note that
$\Gamma(D_{+}(f), \bigoplus_{n \in \mathbf{Z}} \mathcal{O}_X(n)) = A_f$
as a $\mathbf{Z}$-graded algebra. Consider the two
$\mathbf{Z}$-graded ring maps
$$
\theta, \theta' :
A_f
\longrightarrow
\Gamma(T_{\psi(f)}, \bigoplus \mathcal{L}^{\otimes n}).
$$
We know that multiplication by $f$ (resp.\ $\psi(f)$)
is an isomorphism on the left (resp.\ right) hand side.
We also know that $\theta(x/1) = \theta'(x/1) = \psi(x)|_{T_{\psi(f)}}$
by (2) for all $x \in A$. Hence we deduce easily that $\theta = \theta'$
as desired. Considering the degree $0$ parts we deduce that
$r^\sharp = (r')^\sharp$, i.e., that $r = r'$ as morphisms of schemes.
This proves the uniqueness.

\medskip\noindent
Now we come to existence. By the uniqueness just proved, it is enough to
construct the pair $(r, \theta)$ locally on $T$. Hence we may assume
that $T = \Spec(R)$ is affine, that $\mathcal{L} = \mathcal{O}_T$
and that for some $f \in A_{+}$ homogeneous we have
$\psi(f)$ generates $\mathcal{O}_T = \mathcal{O}_T^{\otimes \deg(f)}$.
In other words, $\psi(f) = u \in R^*$ is a unit. In this case the map
$\psi$ is a graded ring map
$$
A \longrightarrow R[x] = \Gamma_*(T, \mathcal{O}_T)
$$
which maps $f$ to $ux^{\deg(f)}$. Clearly this extends (uniquely) to
a $\mathbf{Z}$-graded ring map $\theta : A_f \to R[x, x^{-1}]$ by
mapping $1/f$ to $u^{-1}x^{-\deg(f)}$. This map in degree zero gives
the ring map $A_{(f)} \to R$ which gives the morphism
$r : T = \Spec(R) \to \Spec(A_{(f)}) = D_{+}(f) \subset X$.
Hence we have constructed $(r, \theta)$ in this special case.

\medskip\noindent
Let us show the last statement of the lemma.
According to Lemma \ref{lemma-converse-construction}
the morphism constructed there is the unique one such that
the displayed diagram in its statement commutes.
The commutativity of the diagram in the lemma implies the
commutativity when restricted to $V$ and $A^{(d)}$.
Whence the result.
\end{proof}

\begin{remark}
\label{remark-not-in-invertible-locus}
Assumptions as in Lemma \ref{lemma-invertible-map-into-proj} above.
The image of the morphism $r_{\mathcal{L}, \psi}$ need not be
contained in the locus where the sheaf $\mathcal{O}_X(1)$
is invertible.
Here is an example.
Let $k$ be a field.
Let $S = k[A, B, C]$ graded by $\deg(A) = 1$, $\deg(B) = 2$, $\deg(C) = 3$.
Set $X = \text{Proj}(S)$.
Let $T = \mathbf{P}^2_k = \text{Proj}(k[X_0, X_1, X_2])$.
Recall that $\mathcal{L} = \mathcal{O}_T(1)$ is invertible
and that $\mathcal{O}_T(n) = \mathcal{L}^{\otimes n}$.
Consider the composition $\psi$ of the maps
$$
S \to k[X_0, X_1, X_2] \to \Gamma_*(T, \mathcal{L}).
$$
Here the first map is $A \mapsto X_0$, $B \mapsto X_1^2$,
$C \mapsto X_2^3$ and the second map is (\ref{equation-global-sections}).
By the lemma this corresponds to a morphism
$r_{\mathcal{L}, \psi} : T \to X = \text{Proj}(S)$
which is easily seen to be surjective. On the other hand, in
Remark \ref{remark-not-isomorphism} we showed that the sheaf
$\mathcal{O}_X(1)$ is not invertible at all points of $X$.
\end{remark}










